\begin{abstract}
A staged search must decide which candidate to advance next and when to stop. Some observations
only inform later decisions; others can be kept as results. We formulate these choices as a
staged Pandora's box with a designated set of collectable stages. When each candidate's
readings and positive bounded costs depend only on its own history, an exact covered-call
recursion gives the reservation index, the optimal value and the loss of every advance or stop;
advancing an index-maximal candidate while its index exceeds the incumbent is optimal. The
representation recovers classical single-stage search, nested inspection and the
all-collectable Markov setting, and prices changes in the problem: enlarging the collectable set
or lowering costs raises the value by an explicit covered-call overshoot. Resumption can be
valuable even when future reward laws do not change, but has no value for deterministic
readings. Information sharing limits index optimality: no rule using only the candidates'
separate remaining-reading laws is universally optimal, and making an intermediate reading
collectable can change which candidate is inspected first. Under a coherent joint path law,
with costs that may depend on every candidate's history, the recomputed index's regret is at
most twice the horizon times a uniform bound on how much one reading changes other candidates'
local values, through their laws or costs; this bound is zero under independence with
own-history costs, need not be small under posterior concentration alone, and has a common-component
envelope through predictive total-variation dispersion. Exact finite-law experiments check the
representation, the loss accounting and the bounds.
\end{abstract}

\section{Introduction}
\label{sec:intro}

Cost-aware search has two coupled decisions, in Weitzman's terms a \emph{selection rule},
which candidate to inspect next, and a \emph{stopping rule}, whether to inspect at all
\citep{weitzman1979}. For independent candidates inspected in one step, one reservation value
answers both: open a box of largest reservation value while that value exceeds the best
result already available. The Pandora's box Gittins index (PBGI) brings this principle to
cost-aware Bayesian optimization \citep{xie2024pbgi,xie2025coststop}.

An evaluation need not reveal its result at once. A training run may be paused at checkpoints
and resumed later; a proxy evaluation may inform the decision to pay for a full one without
being a deployable result; a learning curve may make every checkpoint a usable result. We
represent these situations with a finite ladder per candidate and a \emph{collectable set}:
only scores at designated stages enter the incumbent. The searcher may advance any available
candidate by one stage, return to another, or stop and keep its best collectable reading less
the costs incurred. Staging changes both rules: selection includes resuming a paused candidate,
and what stopping keeps depends on the collectable set.

In multi-fidelity Bayesian optimization, a candidate can be a fixed hyperparameter
configuration and its ladder a prescribed sequence of fidelities or training checkpoints.
Terminal-only collection models proxy evaluations whose outputs cannot be deployed;
checkpoint collection models intermediate outputs that are themselves admissible results.
We study the search structure on a finite candidate pool under bounded planning laws;
the numerical studies verify this structure on known laws.

We give an explicit covered-call representation of this problem. At a nonexhausted arm
state $s$ with a collectable stage left, let $\Yc$ be the next stage's score if it
is collectable and $-\infty$ otherwise, and let $s'$ be the successor own state.
The recursion
\[
W(s)=\max\{\Yc,Z(s')\},\qquad Z(s)=\min\{G(s),W(s)\},\qquad
\E[(W(s)-G(s))^+]=c(s),
\]
represents each remaining arm by its next stage's collectable reading or future capped
value, capped again at its current reservation value $G(s)$. A suffix without a collectable
stage, including an exhausted arm, has $G=Z=-\infty$ and needs no reservation equation.
Under independent arms the optimal value is the expectation of the largest
covered call value and the incumbent (Theorem~\ref{thm:index}). The same representation gives
the exact Bellman loss of every advance and stop, hence of any policy
(Theorem~\ref{thm:loss}), and prices changes in the problem: enlarging the collectable set or
lowering costs raises the value by an explicit overshoot, planning with a smaller collectable
set loses at most that amount, and deterministic readings need no resumption.
With stochastic collectable rewards, resumption can help without learning about
unobserved rewards (Section~\ref{sec:comparative}).

We do not claim a new index theorem: the policy is consistent with classical superprocess
results \citep{whittle1980gittins,glazebrook1982whittle}, and the single-stage, terminal-only
and all-collectable Markov cases have established solutions (Section~\ref{sec:related}). Our
contribution is the explicit recursion, value representation and action losses for designated
collection, together with the limits and stability of the index under sharing. The construction separates information from collectable reward
while accommodating own-history states and costs, features already permitted in nested
inspection models \citep[Section~4.1]{bowers2024nested}.

Information sharing marks the limits of these guarantees. Under a shared surrogate, one
candidate's observation changes the law of the others, and separate arm problems ignore that
change, which can raise or lower the value. Two instances with identical separate
arm laws can have different uniquely optimal first actions (Proposition~\ref{prop:no-local}),
and collection rights reshape acquisition: in a two-arm family with a common
latent state, the optimal first advance alternates $A,B,A,B$ in the prior when the second
arm's intermediate reading is collectable, and is always $A$ when it is not
(Proposition~\ref{prop:switching}). These results locate the boundary of index optimality;
they do not propose a coupled planner. Past that boundary the index is
stable when sharing is uniformly small: for a coherent joint path law, its regret over $H$
advances is at most $2H\bar\varepsilon$, where $\bar\varepsilon$ bounds how much one reading
changes the untouched arms' local keep-or-run values, through their laws or costs, at every
feasible future history (Proposition~\ref{prop:sharing-robustness}); $\bar\varepsilon=0$ under
independence with own-history costs.
Through a common component, sensitivity times predictive total-variation dispersion bounds
these changes (Corollary~\ref{cor:component}). Small root covariance or posterior
concentration alone is insufficient, and these conditions are not established for a general
correlated Gaussian process.

\subsection{Contributions}
\label{sec:contributions}
We provide (i) the staged model and a covered-call recursion giving, for arms whose laws and
costs depend only on their own histories, the optimal value and index policy; (ii) exact stop
and advance losses with an arbitrary-policy decomposition, exact comparative statics for
collection rights and costs, and a bound on the value of resumption; (iii) the limits and
stability of the index under sharing, from the own-arm-law and switching results to the regret
bound $2H\bar\varepsilon$ and its common-component envelope; and (iv) exact finite-state checks of
the value and loss identities, comparative statics and sharing bounds against direct Bellman optimization.
